% ------------------------------------------------------------------------------
% Slide 8: Límites de los Sensores Actuales y Candidato para la Lengua
% ------------------------------------------------------------------------------
\begin{frame}{Límites de los Sensores Actuales y Candidato para la Lengua}
    \begin{columns}[c]
        \begin{column}{0.51\textwidth}
            \scriptsize
            \begin{alertblock}{Mezcla de Señales en las Posiciones Actuales}
                \begin{itemize}\setlength{\itemsep}{1pt}\setlength{\parskip}{0pt}
                    \item \textbf{Centro del mentón (digástrico):} Se activa al abrir la boca en la /a/, pero también salta si la lengua sube en la /i/, pisándose dos picos en el mismo canal.
                    \item \textbf{Mejilla (cigomático):} Activo al sonreír en la /i/, pero también se dispara en la /a/ si el hablante abre grande la boca hacia abajo.
                \end{itemize}
            \end{alertblock}
            \vspace{0.1cm}
            \begin{exampleblock}{Nueva Posición: Abajo del Mentón al Costado}
                Colocar un sensor en la \textbf{zona lateral debajo de la mandíbula} (lee la lengua sin el choque central) $+$ sensor de apertura $+$ sensor en labios (\textit{orbicular}). Cubre directamente las 3 variables del modelo de Trevisan y Mindlin.
            \end{exampleblock}
        \end{column}
        \begin{column}{0.47\textwidth}
            \centering
            \placeholderfigura{0.70\textheight}{Fotos / Mapa de ubicación de los electrodos mostrando la competencia de picos y la nueva posición submandibular lateral}
        \end{column}
    \end{columns}
\end{frame}
